\scanpage{47}
\mspage{38}

$\beta)$ Pour tout $i\in I$, $X_i$ est saturé par la relation \deleted{\unsure{d'ordre}} \inserted{d'équivalence} définie par cette partition.

$\gamma)$ Pour $i,j\in I$, \deleted{\unsure{avec $X_j^{\circ}\ne0$}, \unsure{la relation on a}}
\[
X_j^{\circ}\subset X_i\ \inserted{\text{et }X_j^{\circ}\ne\varnothing}\Longrightarrow j\leq\deleted{\unsure{j}}i.
\]
% margin note written slanted
\margin{\uline{Question} $\beta)$ et $\gamma)$ conséquences de $\alpha)$ ? (OK si $I$ fini) non, cf page 44 exemple où on a $\alpha$, \deleted{\unsure{\gamma}} et $\gamma$, mais pas $\beta$. Donc quid de $\alpha$ et $\beta$ ? $\gamma$ conséquence ? Non plus, voir deuxième exemple \unsure{loc. cit.|ibid.}}

\uline{Dém} La condition (i), \deleted{\unsure{i.e. que des $I^x$ signifie}} pour tout \deleted{$\unsure{x\in X}$, \unsure{ce qu'il}} \inserted{$x\in|\Phi|$, $I^x$} ait un plus petit élément, signifie \deleted{trois choses :} $(\forall x\in X)$ trois choses :

$1^{\circ})$ $I^x$ \inserted{(dont on sait qu'il est $\ne\varnothing$)} a un élément minimal, $i$, i.e. $\exists i\in I$ tel que \deleted{$\unsure{I^x}$ est} soit él. min. de $I^x$ i.e.
\[
x\in X_i^{\circ}.
\]

$2^{\circ})$ Deux él. minimaux de $I^x$ sont égaux i.e.
\[
\forall i,j\in I,\quad x\in X_i^{\circ}\cap X_j^{\circ}\Longrightarrow i=j
\]

$3^{\circ})$ Tout élément de $I^x$ est minoré par un élément minimal, \deleted{\unsure{donc}} i.e.
\[
\forall i\in I\text{ tel que }x\in X_i,\quad \exists j\leq i\text{ tel que }x\in X_j^{\circ}.
\]
% margin note written slanted
\margin{\unsure{Noter que} $3^{\circ})\Rightarrow1^{\circ})$ (car on sait $I^x\ne\varnothing$)}
C'est donc la conjonction des trois conditions : $\forall x$ on a $1^{\circ}$, $\forall x$ on a $2^{\circ}$, $\forall x$ on a $3^{\circ}$. Cela s'écrit
\[
|\Phi|=\bigcup X_i^{\circ}\qquad\text{(autom. si }I\text{ fini)}\tag*{$1'$)}
\]
\[
[X_i^{\circ}\cap X_j^{\circ}\ne\varnothing\Rightarrow i=j\quad\text{i.e.}]\quad i\ne j\Longrightarrow X_i^{\circ}\cap X_j^{\circ}=\varnothing\tag*{$2'$)}
\]
\[
\forall i\in I,\text{ on a}\quad X_i=\bigcup_{j\leq i}X_j^{\circ}\quad\text{(autom. si }I\text{ fini)}\tag*{$3'$)}
\]
% margin note written slanted
\margin{\uline{NB} $1',3'$ tjrs \unsure{satisfaites} si $I$ fini. Dans ce cas, la seule condition qui reste est $\forall i,j$, $(i\ne j)$ on a $X_i^{\circ}\cap X_j^{\circ}=\varnothing$.}
% margin note written slanted
\margin{Question : $1'+2'\Longrightarrow3'$ ??? (OK si $I$ fini) non, cf p.~44}
On voit d'ailleurs, puisque $|\Phi|=\bigcup_{i\in I}X_i$, que $3')$ implique $1')$. Donc la condition \unsure{$1^{\circ})$|(i)} équivaut à la conjonction de $2'$ et $3'$, qui ne sont autres que a), b) de (iii). Ainsi
\[
(i)\Longleftrightarrow(iii).
\]

D'autre part, les parties $\alpha)$, $\beta$, $\gamma)$ de (iv) se transcrivent respectivement en :
\[
|\Phi|=\bigcup_{i\in I'}\{X_i^{\circ}\}\quad\text{(équivalent à }1'\text{)}\quad\deleted{\text{et }i,j\in I',}\tag*{$1''$)}
\]
\[
\deleted{2'')\quad\forall i\in I,\text{ posant }I(i)=}
\]
\uline{et}
\[
\forall i,j\in I',\quad\deleted{X_i^{\circ}\cap X_j^{\circ}}\text{ on a }X_i^{\circ}\cap X_j^{\circ}=\varnothing\quad\text{(équivalent à }2'\text{)}.
\]
% archive sheet number: 47
